\section{Invertibility and structure}

It is possible to treat invertibility as a purely computational question: compute the determinant, or try Gauss--Jordan elimination, and see what happens. This works, but it is not the only way to think. Often the structure of a matrix already tells us something about whether an inverse should exist.

The central rule remains the same:
\[
A \text{ is invertible } \iff \det A\ne0.
\]
But before computing the determinant, we should learn to look at the matrix.

\subsection{Diagonal matrices}

A diagonal matrix is one of the clearest cases. Let
\[
D=\operatorname{diag}(d_1,d_2,\ldots,d_n).
\]
Then
\[
\det D=d_1d_2\cdots d_n.
\]
Therefore \(D\) is invertible exactly when none of the diagonal entries is zero.

If every \(d_i\ne0\), then
\[
D^{-1}=\operatorname{diag}\left(\frac1{d_1},\frac1{d_2},\ldots,\frac1{d_n}\right).
\]

\begin{examplebox}
Let
\[
D=\begin{pmatrix}
4&0&0\\
0&-2&0\\
0&0&5
\end{pmatrix}.
\]
Then
\[
D^{-1}=\begin{pmatrix}
\frac14&0&0\\
0&-\frac12&0\\
0&0&\frac15
\end{pmatrix}.
\]
This inverse is easy to understand: each coordinate was scaled independently, so the inverse simply rescales each coordinate back.
\end{examplebox}

If a diagonal entry is zero, one coordinate is collapsed to zero. That information cannot be recovered.

\subsection{Triangular matrices}

Triangular matrices are also easy to test. If \(T\) is triangular, then
\[
\det T=\text{product of diagonal entries}.
\]
Thus \(T\) is invertible if and only if every diagonal entry is non-zero.

\begin{examplebox}
Let
\[
T=\begin{pmatrix}
1&2&3\\
0&-1&4\\
0&0&6
\end{pmatrix}.
\]
Since
\[
\det T=1\cdot(-1)\cdot6=-6\ne0,
\]
the matrix is invertible.

On the other hand,
\[
S=\begin{pmatrix}
2&1&0\\
0&0&5\\
0&0&3
\end{pmatrix}
\]
has determinant
\[
\det S=2\cdot0\cdot3=0.
\]
So \(S\) is not invertible.
\end{examplebox}

This is not merely a shortcut. A zero on the diagonal of a triangular matrix means that one pivot is missing. Without all pivots, the matrix cannot be reduced to the identity matrix.

\subsection{Repeated or dependent rows}

If two rows of a matrix are equal, the determinant is zero. If one row is a multiple of another, the determinant is also zero. More generally, if one row can be built from the others, the determinant is zero.

\begin{examplebox}
Consider
\[
A=\begin{pmatrix}
1&2&1\\
3&6&3\\
0&1&4
\end{pmatrix}.
\]
The second row is three times the first row. Therefore the rows are dependent, and
\[
\det A=0.
\]
So \(A\) is not invertible.
\end{examplebox}

The interpretation is important. If one row repeats information already contained in another row, then the matrix does not contain enough independent information to be reversed.

\subsection{Products of matrices}

The determinant rule
\[
\det(AB)=\det A\det B
\]
has a direct consequence for invertibility.

\begin{theorembox}
If \(A\) and \(B\) are invertible square matrices of the same size, then \(AB\) is invertible, and
\[
(AB)^{-1}=B^{-1}A^{-1}.
\]
\end{theorembox}

The order in this formula is essential. To undo the product \(AB\), we must undo the last action first.

Suppose a vector \(x\) is transformed by \(B\) and then by \(A\):
\[
x \longmapsto Bx \longmapsto A(Bx)=ABx.
\]
To recover \(x\), first undo \(A\), and then undo \(B\):
\[
ABx \longmapsto Bx \longmapsto x.
\]
Thus the inverse is
\[
B^{-1}A^{-1}.
\]

We can verify algebraically:
\[
(AB)(B^{-1}A^{-1})=A(BB^{-1})A^{-1}=AI A^{-1}=AA^{-1}=I.
\]
Similarly,
\[
(B^{-1}A^{-1})(AB)=B^{-1}(A^{-1}A)B=B^{-1}IB=B^{-1}B=I.
\]

\subsection{Matrices satisfying special identities}

Sometimes a matrix is invertible because it satisfies an identity that already reveals its inverse.

\begin{examplebox}
Suppose a square matrix \(A\) satisfies
\[
A^2=I.
\]
Then \(A\) is its own inverse, because
\[
AA=I.
\]
Therefore
\[
A^{-1}=A.
\]
Such a matrix is called an involution.
\end{examplebox}

\begin{examplebox}
Suppose a square matrix \(A\) satisfies
\[
A^2-3A+2I=0.
\]
Then
\[
A^2-3A=-2I.
\]
Factor the left side:
\[
A(A-3I)=-2I.
\]
Multiplying by \(-\frac12\), we get
\[
A\left(-\frac12(A-3I)\right)=I.
\]
Therefore
\[
A^{-1}=-\frac12(A-3I)=\frac12(3I-A).
\]
This argument finds the inverse without performing Gauss--Jordan elimination.
\end{examplebox}

These examples show that invertibility is not only a computational topic. It is also a structural topic. Sometimes the matrix tells us how it can be undone.

\begin{summarybox}
To recognize invertibility, look for structure:
\begin{itemize}
\item diagonal matrices are invertible when no diagonal entry is zero,
\item triangular matrices are invertible when all diagonal entries are non-zero,
\item repeated or dependent rows imply non-invertibility,
\item products of invertible matrices are invertible,
\item special identities may reveal the inverse directly.
\end{itemize}
\end{summarybox}
